\SourcePage{383}{386}
\section[Singlet-term vibration and rotation]{Vibrational and rotational structure of singlet terms of a diatomic molecule}

As noted at the beginning of this chapter, the large disparity between the nuclear and electronic masses allows the determination of molecular energy levels to be divided into two parts. First one finds, as functions of the internuclear distance, the energy levels of the electrons for fixed nuclei (the electronic terms). One may then consider nuclear motion in a specified electronic state; the nuclei are thereby treated as particles interacting through $U_n(r)$, where $U_n$ is the corresponding electronic term. Molecular motion consists of translation of the molecule as a whole and motion of the nuclei relative to their center of mass. The translational part is of no interest here, and the center of mass may be taken as fixed.

For clarity, consider first electronic terms for which the total molecular spin $S$ vanishes (singlet terms). The relative motion of two particles (the nuclei) interacting through $U(r)$ reduces to the motion of one particle of mass $M$ (their reduced mass) in the central field $U(r)$. (Here $U(r)$ denotes the energy of the electronic term under consideration.) Motion in the central field, in turn, reduces to a one-dimensional problem with an effective energy equal to $U(r)$ plus the centrifugal energy.

Let $\mathbf K$ be the total angular momentum of the molecule, composed of the electronic orbital angular momentum $\mathbf L$ and the angular momentum of nuclear rotation. The operator of the nuclear centrifugal energy is then
\[B(r)(\mathbf K-\mathbf L)^2,\]
where the notation
\begin{equation}B(r)=\frac{\hbar^2}{2Mr^2},\tag{82.1}\end{equation}
customary in the theory of diatomic molecules, has been introduced. Averaging this quantity over the electronic state (at fixed $r$) gives the centrifugal energy as a function of $r$, which must enter the effective potential energy $U_K(r)$. Hence
\begin{equation}U_K(r)=U(r)+B(r)\overline{(\mathbf K-\mathbf L)^2},\tag{82.2}\end{equation}
where the bar denotes this averaging.

\SourcePage{384}{387}
Carry out the average for a state with a definite square of the total angular momentum, $\mathbf K^2=K(K+1)$ ($K$ is an integer), and a definite projection of the electronic angular momentum on the molecular axis (the $z$ axis), $L_z=\Lambda$. Expansion of (82.2) gives
\begin{equation}U_K(r)=U(r)+B(r)K(K+1)-2B(r)\overline{\mathbf{LK}}+B(r)\overline{\mathbf L^2}.\tag{82.3}\end{equation}
The last term depends only on the electronic state and contains no quantum number $K$; it can simply be absorbed into $U(r)$. We now show that this also applies to the penultimate term.

Recall that when the projection of an angular momentum on an axis has a definite value, the mean angular-momentum vector is directed along that same axis (see the remark at the end of \S~27). Let $\mathbf n$ denote the unit vector along the $z$ axis: $\overline{\mathbf L}=\Lambda\mathbf n$. In classical mechanics the rotational angular momentum of a two-particle system (the nuclei) is $\mathbf r\times\mathbf p$, where $\mathbf r=r\mathbf n$ is the radius vector between the particles and $\mathbf p$ is the momentum of their relative motion; this quantity is perpendicular to $\mathbf n$. The same statement holds in quantum mechanics for the operator of nuclear rotational angular momentum: $(\mathbf K-\mathbf L)\mathbf n=0$, or $\mathbf K\mathbf n=\mathbf L\mathbf n$. Equality of the operators entails equality of their eigenvalues, and since $\mathbf n\mathbf L=L_z=\Lambda$, it follows that
\begin{equation}\mathbf K\mathbf n=\Lambda.\tag{82.4}\end{equation}
Thus, in the penultimate term of (82.3), $\overline{\mathbf{LK}}=\mathbf n\mathbf K\Lambda=\Lambda^2$, which is independent of $K$. After redefining $U(r)$, the effective potential energy finally becomes
\begin{equation}U_K(r)=U(r)+B(r)K(K+1).\tag{82.5}\end{equation}
We also note that $K_z=\Lambda$ implies that, for a given $\Lambda$, the quantum number $K$ can assume only the values
\begin{equation}K\geqslant\Lambda.\tag{82.6}\end{equation}

Solving the one-dimensional Schrödinger equation with the potential energy (82.5) yields a sequence of energy levels. For each fixed $K$, let these levels be numbered in ascending order by $v=0,1,2,\ldots$, with $v=0$ corresponding to the lowest one. Nuclear motion therefore splits every electronic term into a set of levels characterized by the two quantum numbers $K$ and $v$.

\SourcePage{385}{388}
For a given electronic term, the number of these levels may be finite or infinite. Suppose the electronic state is such that, as $r\to\infty$, the molecule becomes two separate neutral atoms. Then $U(r)$, and with it $U_K(r)$, tends to the limiting constant $U(\infty)$ (the sum of the energies of the two isolated atoms) faster than $1/r^2$ (see \S~89). Such a field has finitely many levels (see \S~18), although their actual number in molecules is very large. Their distribution is as follows: for each fixed $K$ there is a certain number of levels distinguished by $v$, and this number decreases as $K$ increases, until a value of $K$ is reached for which no levels remain at all.

If, however, the molecule separates into two ions as $r\to\infty$, then at large distances $U(r)-U(\infty)$ approaches the Coulomb attraction energy of the ions ($\sim1/r$). Such a field has infinitely many levels, which crowd together on nearing the limiting value $U(\infty)$. We note that the first case applies to most molecules in the normal state; only a comparatively small number of molecules yield ions when their nuclei are drawn apart.


A complete general calculation of how the energy levels depend on the quantum numbers is not possible. It can be carried out only for levels with relatively low excitation, whose energies do not lie far above the ground level\footnote{All levels considered here originate from one specified electronic term.}. For such levels the quantum numbers $K$ and $v$ are small. Since molecular spectroscopy normally deals with levels of this kind, they are of particular interest.


In weakly excited states, nuclear motion can be described as small oscillations about equilibrium. Accordingly, expand $U(r)$ in powers of $\xi=r-r_e$, where $r_e$ is the value of $r$ at which $U(r)$ is minimal. Since $U'(r_e)=0$, through second order $U(r)=U_e+M\omega_e^2\xi^2/2$, where $U_e=U(r_e)$ and $\omega_e$ is the oscillation frequency. In the second term of (82.5), the centrifugal energy, it is sufficient to put $r=r_e$, since this term already contains the small quantity $K(K+1)$. We thus have
\begin{equation}U_K(r)=U_e+B_eK(K+1)+\frac{M\omega_e^2\xi^2}{2},\tag{82.7}\end{equation}

\SourcePage{386}{389}
where $B_e=\hbar^2/(2Mr_e^2)=\hbar^2/2I$ is the so-called \emph{rotational constant} ($I=Mr_e^2$ is the molecule's moment of inertia).


The first two terms in (82.7) are constants, whereas the third represents a one-dimensional harmonic oscillator. The energy levels being sought may therefore be written down immediately as

\begin{equation}E=U_e+B_eK(K+1)+\hbar\omega_e\left(v+\frac12\right).\tag{82.8}\end{equation}
Thus, in this approximation an energy level consists of three independent parts:
\begin{equation}E=E^{el}+E^r+E^v.\tag{82.9}\end{equation}
The first term ($E^{el}=U_e$) is the electronic energy (including the energy of the Coulomb interaction between the nuclei at $r=r_e$). The second term

\begin{equation}E^r=B_eK(K+1),\tag{82.10}\end{equation}
~--- the \emph{energy of rotation} (or rotational energy) associated with the molecule's rotation\footnote{The wave function describing the rotation of a diatomic molecule (with spin disregarded) has essentially the same form as the wave function of a symmetric top (\S~103). Unlike the top, a molecule's rotation is described by only two angles ($\alpha\equiv\varphi$, $\beta\equiv\theta$), which specify the direction of its axis. The rotational wave function differs from (103.8) in lacking the factor $e^{ik\gamma}/\sqrt{2\pi}$ and in the notation used for the quantum numbers. Since, by (82.4), $\Lambda$ equals the projection of the total angular momentum $\mathbf K$ on the molecular axis (the $\zeta$ axis in \S~103), the notation must be changed according to $J,M,k\to K,M,\Lambda$ (where now $M=K_z$). Thus, $\psi_{\mathrm{rot}}(\varphi,\theta)=i^K\sqrt{(2K+1)/(4\pi)}D_{\Lambda M}^{(K)}(\varphi,\theta,0)$.}. Finally, the third term
%
\unskip,
\begin{equation}E^v=\hbar\omega_e\left(v+\frac12\right),\tag{82.11}\end{equation}
is the energy of nuclear vibration within the molecule. In accordance with its definition, $v$ numbers in ascending order the levels having a given $K$; it is called the \emph{vibrational} quantum number.

For a fixed shape of the potential-energy curve $U(r)$, the frequency $\omega_e$ is inversely proportional to $\sqrt M$. The spacings $\Delta E^v$ between vibrational levels are consequently proportional to $1/\sqrt M$. The spacings $\Delta E^r$ between rotational levels contain the moment of inertia $I$ in the denominator and are therefore proportional to $1/M$. By contrast, the spacings $\Delta E^{el}$ between electronic levels,

\SourcePage{387}{390}
like the levels themselves, contain no $M$. Since $m/M$ ($m$ being the electron mass) is the small parameter in the theory of diatomic molecules, we find
\begin{equation}\Delta E^{el}\gg\Delta E^v\gg\Delta E^r.\tag{82.12}\end{equation}
These inequalities express the distinctive distribution of molecular energy levels. Nuclear vibration resolves the electronic terms into levels lying comparatively close together. Molecular rotation, in turn, produces a still finer splitting of these levels\footnote{As an example, values of $U_e$, $\hbar\omega_e$, and $B_e$ (in electron-volts) for several molecules are
\[\begin{array}{c|ccc}&\mathrm H_2&\mathrm N_2&\mathrm O_2\\\hline-U_e&4{,}7&7{,}5&5{,}2\\\hbar\omega_e&0{,}54&0{,}29&0{,}20\\10^3B_e&7{,}6&0{,}25&0{,}18\end{array}\] }.

In subsequent approximations the energy can no longer be separated into independent vibrational and rotational parts; rotation--vibration terms appear that contain both $K$ and $v$. Successive approximations would give the levels $E$ as a power-series expansion in the quantum numbers $K$ and $v$.

We calculate here the approximation immediately following (82.8). For this purpose the expansion of $U(r)$ in powers of $\xi$ must be continued through fourth-order terms (compare the anharmonic-oscillator problem in \S~38). Correspondingly, expand the centrifugal energy through terms in $\xi^2$. This gives
\begin{equation}\begin{split}U_K(r)={}&U_e+\frac{M\omega_e^2}{2}\xi^2+\frac{\hbar^2}{2Mr_e^2}K(K+1)-a\xi^3+b\xi^4\\&-\frac{\hbar^2}{Mr_e^3}K(K+1)\xi+\frac{3\hbar^2}{2Mr_e^4}K(K+1)\xi^2.\end{split}\tag{82.13}\end{equation}
Now calculate the correction to the eigenvalues (82.8), treating the last four terms in (82.13) as the perturbation operator. First-order perturbation theory suffices for the terms in $\xi^2$ and $\xi^4$, whereas second order must be calculated for those in $\xi$ and $\xi^3$, since the diagonal matrix elements of $\xi$ and $\xi^3$ vanish identically. All matrix elements needed for the calculation were found in \S~23 and in problem 3 of \S~38. The result is

\SourcePage{388}{391}
conventionally written as
\begin{equation}\begin{split}E={}&E^{el}+\hbar\omega_e\left(v+\frac12\right)-x_e\hbar\omega_e\left(v+\frac12\right)^2\\&+B_vK(K+1)-D_eK^2(K+1)^2,\end{split}\tag{82.14}\end{equation}
where
\begin{equation}B_v=B_e-\alpha_e\left(v+\frac12\right)\equiv B_0-\alpha_ev.\tag{82.15}\end{equation}
The constants $x_e$, $B_e$, $\alpha_e$, and $D_e$ are related to the constants entering (82.14) by
\begin{equation}\begin{gathered}B_e=\frac{\hbar^2}{2I},\qquad D_e=\frac{4B_e^3}{\hbar^2\omega_e^2},\\
\alpha_e=\frac{6B_e^2}{\hbar\omega_e}\left(\frac{a\hbar}{M\omega_e^2}\sqrt{\frac2{MB_e}}-1\right),\\
x_e=\frac3{2\hbar\omega_e}\left(\frac{\hbar}{M\omega_e}\right)^2\left(\frac52\frac{a^2}{M\omega_e^2}-b\right).
\end{gathered}\tag{82.16}\end{equation}
Terms independent of $v$ and $K$ have been included in $E^{el}$.

\ProblemHeading Estimate the accuracy of the approximation that separates electronic and nuclear motion in a diatomic molecule.

\SolutionHeading Write the complete molecular Hamiltonian as $\widehat H=\widehat T_r+\widehat H_{el}$, where $\widehat T_r=\widehat{\mathbf p}^{2}/2M$ is the kinetic-energy operator for the relative motion of the nuclei ($\widehat{\mathbf p}=-i\hbar\partial/\partial\mathbf r$; $\mathbf r$ is the vector separating the nuclei, and $M$ their reduced mass). The Hamiltonian $\widehat H_{el}$ includes the electron kinetic-energy operators, the potential energy of their Coulomb interactions with one another and with the nuclei, and the Coulomb interaction energy of the nuclei\footnote{The Hamiltonian $\widehat H$ refers to a frame in which the center of mass of the whole molecule is at rest ($\mathbf P_n+\mathbf P_e=0$, where $\mathbf P_n$ is the total momentum of the two nuclei and $\mathbf P_e$ the total electronic momentum). The term corresponding to the kinetic energy of the nuclear center of mass has nevertheless already been omitted: $\widehat{\mathbf P}_n^{2}/2(M_1+M_2)=\mathbf p_e^2/2(M_1+M_2)$. This term is certainly smaller than the electronic kinetic energy by the ratio $m/M$.}. Seek a solution of the Schrödinger equation
\begin{equation}\widehat H\psi=(\widehat T_r+\widehat H_{el})\psi=E\psi\tag{1}\end{equation}
in the form
\begin{equation}\psi=\sum_m\chi_m(\mathbf r)\varphi_m(q,\mathbf r),\tag{2}\end{equation}

\SourcePage{389}{392}
where the functions $\varphi_m(q,\mathbf r)$ are orthonormal solutions of
\begin{equation}\widehat H_{el}\varphi_m(q,\mathbf r)=U_m(r)\varphi_m(q,\mathbf r)\tag{3}\end{equation}
($q$ denotes the full set of electronic coordinates), and $U_m(r)$ are the eigenvalues of $\widehat H_{el}$ with $r$ treated as a parameter. Substitute (2) into (1), multiply on the left by $\varphi_n^*(q,\mathbf r)$, and integrate over $dq$. The result is
\begin{equation}\left[\frac{\widehat{\mathbf p}^{2}}{2M}+V_{nn}''+U_n(r)-E\right]\chi_n(\mathbf r)=-\sum_m'(\widehat V_{nm}'+V_{nm}'')\chi_m(\mathbf r),\tag{4}\end{equation}
where
\[\widehat V_{nm}'=\frac1M\mathbf p_{nm}\widehat{\mathbf p},\qquad V_{nm}''=\frac1{2M}(\mathbf p^2)_{nm},\]
and $\mathbf p_{nm}=\int\varphi_n^*\widehat{\mathbf p}\varphi_m,dq$, while $(\mathbf p^2)_{nm}$ is likewise a matrix element with respect to the electronic wave functions; the diagonal element $\mathbf p_{nn}$ vanishes by symmetry.

The electronic functions $\varphi_n$ vary appreciably only over distances of atomic order. Differentiating them with respect to $\mathbf r$ therefore introduces no large parameter $M/m$ ($m$ is the electron mass). Consequently, $V_{nn}''$ is smaller than $U_n(r)$ by the ratio $m/M$ and may be omitted. If the terms on the right-hand side of (4) are regarded as a small perturbation, then to zeroth order the functions $\chi_n(\mathbf R)$ are solutions of
\begin{equation}\left[\frac{\widehat{\mathbf p}^{2}}{2M}+U_n(r)\right]\chi_{nv}=E_{nv}\chi_{nv},\tag{5}\end{equation}
which describes nuclear motion in the field $U_n(r)$ ($v$ denotes the quantum numbers of this motion). Perturbation theory is applicable provided
\[|\langle nv'|\widehat V_{nm}'+V_{nm}''|mv\rangle|\ll|E_{nv'}-E_{mv}|.\]
On the right are energy differences belonging to different electronic terms; they are of zeroth order in the small parameter $m/M$. On the left are matrix elements over the nuclear wave functions. The term involving $\widehat V_{nm}''$ contains $m/M$ and is certainly small. In the matrix element of $\widehat V_{nm}'$, the operator $\widehat{\mathbf p}$ acting on $\chi_{mv}$ multiplies it by a quantity of the order of the nuclear momentum. For small nuclear oscillations this momentum is $\sim\sqrt{M\hbar\omega_e}$; since at the same time $\omega_e$ is inversely proportional to $\sqrt M$, the matrix element $\langle nv'|\widehat V_{nm}'|mv\rangle$ is of order $(m/M)^{3/4}$.
